\documentclass[11pt]{article}
\usepackage{amsmath, amsthm, amssymb}
\usepackage{mathtools}
\usepackage{enumitem}
\usepackage[margin=1in]{geometry}
\usepackage{booktabs}
\usepackage{hyperref}

\theoremstyle{plain}
\newtheorem{theorem}{Theorem}
\newtheorem{lemma}[theorem]{Lemma}
\newtheorem{proposition}[theorem]{Proposition}
\newtheorem{corollary}[theorem]{Corollary}
\newtheorem{conjecture}[theorem]{Conjecture}

\theoremstyle{definition}
\newtheorem{definition}[theorem]{Definition}
\newtheorem{example}[theorem]{Example}

\theoremstyle{remark}
\newtheorem{remark}[theorem]{Remark}
\newtheorem{observation}[theorem]{Observation}

\newcommand{\Z}{\mathbb{Z}}
\newcommand{\Q}{\mathbb{Q}}
\newcommand{\Sym}{\mathfrak{S}}
\newcommand{\Hq}{H_q}
\newcommand{\Pibs}{\Pi_q}
\newcommand{\Vlam}{V_\lambda}
\newcommand{\trace}{\operatorname{tr}}
\newcommand{\rank}{\operatorname{rank}}

\title{$\sigma_2$-G1 and $\sigma_3$-G1 at the rank-3 frontier $\lambda = (4,2)$, $n=6$}
\author{Clio}
\date{26 April 2026}

\begin{document}
\maketitle

\begin{abstract}
For the Hecke algebra $\Hq(\Sym_6)$ and the staircase Bott--Samelson operator
$\Pibs^{(6)} = \prod_{j=1}^{15}(1 + T_{i_j})$ along the descending-staircase
reduced word for $w_0 \in \Sym_6$, we establish $\sigma_k$-G1 for $k = 1, 2, 3$
on the smallest rank-3 Specht module $V_{(4,2)}$ (dimension 9, with
$\dim V_{(4,2)}^{H_6} = 3$). Explicitly,
\begin{align*}
\sigma_1(V_{(4,2)}) &= q^{2}(1+q)^{3}\, P_1(q), \\
\sigma_2(V_{(4,2)}) &= q^{5}(1+q)^{10}(q^2+4q+1)\, P_{2}^{\text{inner}}(q), \\
\sigma_3(V_{(4,2)}) &= q^{9}(1+q)^{21}(q^2+4q+1)^{3},
\end{align*}
where $P_1$ and $P_2^{\text{inner}}$ are explicit palindromic polynomials in
$\Z_{\geq 0}[q]$ of degree 8. The most striking feature is the perfect cube
$(q^2+4q+1)^3$ appearing in $\sigma_3$: the universal Eulerian-Hecke atom
$q^2+4q+1 = (1,4,1)$, which controls $\sigma_{\text{top}}$ for every rank-2
shape at $n=5$ and $V_{(3,1)}$ at $n=4$, reappears as a cube precisely at the
rank-3 frontier $V_{(4,2)}$, $n=6$.

The proof of the closed forms is an explicit symbolic computation in the
Hoefsmit seminormal basis. Once the closed forms are established, $\sigma_k$-G1
for $k = 1, 2, 3$ follows by inspection of coefficients. We then unpack the
structural mechanism: $\sigma_3 = \det(\Pibs^{(6)} \mid \mathrm{im})$ where
$\mathrm{im} = V_{(4,2)}^{H_6}$ is 3-dimensional by the H-invariant theorem
(\cite{H-invariant}), and the 3$\times$3 image matrix has cyclotomic
denominators $\Phi_3, \Phi_4$ in its entries which cancel in
$\sigma_2(\mathrm{X})$ and $\sigma_3(\mathrm{X}) = \det(\mathrm{X})$ by
Hecke-character integrality.

We highlight three structurally novel features at this rank-3 frontier:
(i) the cube $(q^2+4q+1)^3$, (ii) two new degree-8 palindromic atoms $P_1$
and $P_2^{\text{inner}}$ outside the previously catalogued 23-atom alphabet,
and (iii) two-sided leak in the $S_5$-branching block decomposition
(rank 1 in both off-diagonal blocks). The structural origin of the cube
remains conjectural; we record it as a precise open question.
\end{abstract}

\section{Setup}

Let $\Hq = \Hq(\Sym_6)$ be the Iwahori--Hecke algebra over $\Q(q)$ with
generators $T_1,\ldots,T_5$ satisfying $T_i^2 = (q-1)T_i + q$ and the braid
relations $T_i T_{i+1} T_i = T_{i+1} T_i T_{i+1}$, $T_i T_j = T_j T_i$ for
$|i - j| \geq 2$. We use the Hoefsmit seminormal basis: for any partition
$\lambda \vdash n$, the Specht module $V_\lambda$ has basis indexed by
standard Young tableaux of shape $\lambda$, with $T_i$ acting by explicit
2$\times$2 blocks indexed by pairs of tableaux differing by the swap of $i$
and $i+1$.

For $\lambda = (4,2)$ at $n = 6$, $\dim V_{(4,2)} = 9$ and the generators
$T_1, \ldots, T_5$ are explicit $9 \times 9$ matrices over $\Q(q)$ satisfying
the Hecke relations (verified symbolically).

Let
\[
\underline{w}_0 = (1;\; 2,1;\; 3,2,1;\; 4,3,2,1;\; 5,4,3,2,1)
\]
be the descending-staircase reduced word for $w_0 \in \Sym_6$ (length 15),
and define
\begin{equation}\label{eq:Pi-def}
\Pibs^{(6)} \coloneqq (1+T_1)(1+T_2)(1+T_1)(1+T_3)(1+T_2)(1+T_1)\cdots(1+T_5)(1+T_4)(1+T_3)(1+T_2)(1+T_1).
\end{equation}
Since $(1+T_i)$ satisfy the same braid relations as $T_i$, the product
$\Pibs^{(6)}$ is independent of the choice of reduced word for $w_0$.

For $k \geq 0$, write
\[
\sigma_k(\Vlam) \coloneqq e_k\bigl(\text{eigenvalues of } \Pibs^{(6)} \text{ on } \Vlam\bigr) = \trace\bigl(\Lambda^k \Pibs^{(6)} \,\bigm|\, \Lambda^k \Vlam\bigr),
\]
the $k$-th elementary symmetric function of the eigenvalues of $\Pibs^{(6)}$
acting on $V_\lambda$. Equivalently, $\sigma_k$ is the sum of all
$k \times k$ principal minors of any matrix representation of
$\Pibs^{(6)} \mid V_\lambda$.

\begin{definition}[$\sigma_k$-G1]
Let $\sigma_k(\Vlam) \in \Z[q]$. Write $\sigma_k(\Vlam) = q^{a}(1+q)^{b} P(q)$
with $(a,b) \in \Z_{\geq 0}^2$ maximal so that $P(q) \in \Z[q]$ and
$P(0) \neq 0$, $P(-1) \neq 0$. We say $\sigma_k$-G1 holds at
$(\lambda, n)$ if $P(q) \in \Z_{\geq 0}[q]$.
\end{definition}

The companion property $\sigma_k$-G2 (palindromicity of $P$) is automatic
on every Specht module, by the standard palindromy of the Hecke algebra
under $v \mapsto v^{-1}$ ($v = q^{1/2}$); it is verified in the
factorizations below.

\section{H-invariants and rank}

For $n$ even, let $H_n = \langle s_1, s_3, s_5, \ldots, s_{n-1}\rangle \cong
\Sym_2 \times \Sym_2 \times \cdots \times \Sym_2 \subset \Sym_n$ (the
``Young'' parabolic of long-cycle type $2^{n/2}$). Let $V_\lambda^{H_n}$
denote the simultaneous $+q$ eigenspace of $T_{s_1}, T_{s_3}, \ldots, T_{s_{n-1}}$
on $V_\lambda$.

\begin{theorem}[H-invariant theorem; Clio, Apr 17 2026]\label{thm:H-inv}
For all $\lambda \vdash n$,
\[
\rank\bigl(\Pibs^{(n)} \mid V_\lambda\bigr) = \dim V_\lambda^{H_n}.
\]
Equivalently, the image of $\Pibs^{(n)}$ on $V_\lambda$ coincides with
$V_\lambda^{H_n}$, and all eigenvalues of $\Pibs^{(n)}$ on $V_\lambda$ are
strictly positive on this image (with eigenvalue $0$ on the complement).
\end{theorem}

\begin{lemma}\label{lem:rank3}
For $\lambda = (4,2)$, $n = 6$:
\[
\dim V_{(4,2)}^{H_6} = 3,\qquad \rank\bigl(\Pibs^{(6)} \mid V_{(4,2)}\bigr) = 3.
\]
In particular $\sigma_4(V_{(4,2)}) = \sigma_5(V_{(4,2)}) = \cdots = 0$, and
$\sigma_3 = \det(\Pibs^{(6)} \mid V_{(4,2)}^{H_6})$ in any basis of the image.
\end{lemma}

\begin{proof}
The first equality is the H-invariant theorem applied to $(\lambda, n) = ((4,2), 6)$.
For $\dim V_{(4,2)}^{H_6}$, observe that $H_6 = \langle s_1, s_3, s_5\rangle$
acts on $V_{(4,2)}$ with character a sum of products of characters of
$\Sym_2 \times \Sym_2 \times \Sym_2$. A direct character computation
(or symbolic rank computation of the projector
$e_1 e_3 e_5 = \prod_{i \in \{1,3,5\}} (1+T_i)/(1+q)$) gives 3.

Direct symbolic computation of $\rank\Pibs^{(6)}$ on $V_{(4,2)}$ over
$\Q(q)$ confirms rank 3.
\end{proof}

\section{Main theorem}

\begin{theorem}[$\sigma_k$-G1 at the rank-3 frontier]\label{thm:main}
For $\lambda = (4,2)$, $n = 6$, the elementary symmetric functions
$\sigma_k = \sigma_k(\Pibs^{(6)} \mid V_{(4,2)})$ for $k = 1, 2, 3$ admit
the following closed forms:
\begin{align}
\sigma_1 &= q^{2}(1+q)^{3}\, P_1(q), \label{eq:sigma1}\\
\sigma_2 &= q^{5}(1+q)^{10}\, (q^2+4q+1)\, P_2^{\mathrm{inner}}(q),\label{eq:sigma2}\\
\sigma_3 &= q^{9}(1+q)^{21}\, (q^2+4q+1)^{3},\label{eq:sigma3}
\end{align}
where $P_1, P_2^{\mathrm{inner}} \in \Z_{\geq 0}[q]$ are the explicit
palindromic polynomials of degree 8:
\begin{align*}
P_1(q) &= 1 + 13q + 79q^2 + 245q^3 + 365q^4 + 245q^5 + 79q^6 + 13q^7 + q^8,\\
P_2^{\mathrm{inner}}(q) &= 1 + 17q + 103q^2 + 313q^3 + 473q^4 + 313q^5 + 103q^6 + 17q^7 + q^8.
\end{align*}
In particular, $\sigma_k$-G1 holds at $(\lambda, n) = ((4,2), 6)$ for every
$k \geq 0$ (with $\sigma_0 = 1$, $\sigma_k = 0$ for $k \geq 4$).
\end{theorem}

\begin{proof}
The statement is verified by explicit symbolic computation in the Hoefsmit
seminormal basis (Section~\ref{sec:computation}). Once the closed forms
\eqref{eq:sigma1}--\eqref{eq:sigma3} are established as polynomial identities
in $\Z[q]$, the G1 conclusion is a direct inspection: $P_1$ and
$P_2^{\mathrm{inner}}$ have non-negative integer coefficients, and
$(q^2+4q+1) \in \Z_{\geq 0}[q]$. A product of polynomials in $\Z_{\geq 0}[q]$
lies in $\Z_{\geq 0}[q]$, hence each $\sigma_k$ is in $\Z_{\geq 0}[q]$.
\end{proof}

\section{Computational proof of the closed forms}\label{sec:computation}

\subsection{Hoefsmit matrices}

The 9 standard Young tableaux of shape $(4,2)$ are enumerated in row-reading
order. In this basis, each $T_i$ ($i = 1, \ldots, 5$) is an explicit
$9 \times 9$ matrix over $\Q(q)$ given by the Hoefsmit formulas:
\begin{itemize}
\item Diagonal entries: if $i, i+1$ in tableau $T$ are in the same row, the
diagonal entry is $q$; if same column, $-1$; otherwise
$(q-1)/(1 - q^{-d})$ where $d$ is the axial-distance increment.
\item Off-diagonal $2 \times 2$ block: the swap $T \leftrightarrow s_i \cdot T$
contributes the standard Hoefsmit pair, with a $1$ on the row of $s_i \cdot T$
and entry $\alpha_T \alpha_{s_i T} + q$ on the row of $T$.
\end{itemize}
The Hecke relations $(T_i)^2 = (q-1)T_i + q$, $T_i T_{i+1} T_i = T_{i+1} T_i T_{i+1}$,
and commutation $T_i T_j = T_j T_i$ for $|i - j| \geq 2$ are verified
symbolically.

\subsection{Computation of $\Pibs^{(6)}$}

The product $\Pibs^{(6)} \in M_9(\Q(q))$ is computed as a length-15 product
of the 9$\times$9 matrices $(1+T_{i_j})$, in symbolic form (each entry a
rational function in $q$, automatically simplified after each multiplication).

\subsection{Newton identities for $\sigma_k$}

Let $p_j = \trace(\Pibs^{(6)})^j$ for $j \geq 1$. The elementary
symmetric functions $\sigma_k = e_k(\text{eigenvalues})$ are related to the
power sums by Newton's identities:
\[
k \cdot \sigma_k = \sum_{j=1}^{k}(-1)^{j-1}\, \sigma_{k-j}\, p_j,
\]
with $\sigma_0 = 1$. Computing $p_1, p_2, p_3, p_4$ (traces of
$\Pibs^{(6)}, (\Pibs^{(6)})^2, (\Pibs^{(6)})^3, (\Pibs^{(6)})^4$) and
applying Newton's identities yields exact expressions for $\sigma_1, \sigma_2, \sigma_3, \sigma_4$
as elements of $\Q(q)$. We find $\sigma_4 \equiv 0$, confirming
Lemma~\ref{lem:rank3}, and the closed forms \eqref{eq:sigma1}--\eqref{eq:sigma3}
are obtained by symbolic factorization (extracting $q^a$, then dividing by
$(1+q)^b$ until division leaves a remainder).

\subsection{Verification via the $3 \times 3$ image matrix}\label{ssec:image-matrix}

Independently of the trace-based computation, $\sigma_3 = \det(\Pibs^{(6)} \mid V_{(4,2)}^{H_6})$
can be computed as a $3 \times 3$ determinant. Choose three columns
$\Pibs^{(6)} e_{j_1}, \Pibs^{(6)} e_{j_2}, \Pibs^{(6)} e_{j_3}$ of the matrix
$\Pibs^{(6)}$ such that they form a basis of the image (always possible since
$\rank = 3$). Assemble them as a $9 \times 3$ matrix $M$ of rank 3, and
solve $\Pibs^{(6)} M = M X$ for $X \in M_3(\Q(q))$ using the linear system
$M^\top M \cdot X = M^\top \Pibs^{(6)} M$ (or by left-pseudoinverse).

In the Hoefsmit basis above, taking $j_1 = 0, j_2 = 1, j_3 = 2$ (the first
three columns) yields a $3 \times 3$ matrix $X$ whose entries are rational
functions in $q$ with non-trivial denominators. The cyclotomic
factors $\Phi_3 = q^2 + q + 1$ and $\Phi_4 = q^2 + 1$ appear in 8 of the 9
denominators (only $X_{2,0}$ has trivial denominator). Specifically:
\[
X_{0,0} = \frac{q^2(1+q)^7(q^6 + 7q^5 + 21q^4 + 29q^3 + 21q^2 + 7q + 1)}{q^2 + 1},
\]
and similar expressions for the remaining entries (computed symbolically).

\begin{lemma}[Pole cancellation]\label{lem:pole-cancellation}
Despite the cyclotomic denominators $\Phi_3, \Phi_4$ appearing in the
individual entries of the image matrix $X$, the elementary symmetric
functions $\sigma_k(X)$ for $k = 1, 2, 3$ are polynomials in $q$. Equivalently,
the cyclotomic denominators cancel after summing over principal minors.
\end{lemma}

\begin{proof}
$\sigma_k(\Pibs^{(6)} \mid V_{(4,2)}) = \sigma_k(X)$ for $k \leq 3$ since
$\sigma_k$ is invariant under change of basis and $X$ is the restriction
of $\Pibs^{(6)}$ to its (3-dimensional) image. On the other hand,
$\sigma_k(\Pibs^{(6)} \mid V_{(4,2)})$ is also computable directly from
traces of powers of the $9 \times 9$ matrix $\Pibs^{(6)}$, whose entries
are themselves elements of $\Q(q)$ but whose traces lie in $\Z[q]$ (by
integrality of the Hecke character on the cell module $V_\lambda$). Hence
$\sigma_k(X) \in \Z[q]$, and the cyclotomic denominators must cancel.
\end{proof}

\begin{remark}
This is the same Hecke-character-integrality cancellation observed at the
rank-2 shapes $\lambda \in \{(5,1), (3,2,1)\}$ for $\sigma_2$ at $n = 6$
(see \cite{2026-04-25-sigma2-G1-n6}, also \cite{2026-04-26-phi4-artifact}):
the $\Phi_d$ poles in the basis-by-basis matrix entries are basis artifacts,
forced to cancel by integrality of the Hecke character. They are
\emph{not} a sign of parity-purity (Sandvik) or of a Goertzen--Williamson
optimization phenomenon; they are automatic.
\end{remark}

\subsection{Direct computation of $\det(X)$}

Computing $\det(X)$ symbolically as a $3 \times 3$ determinant from the
explicit entries gives
\begin{equation}
\det(X) = q^{9}(1+q)^{21}(q^2+4q+1)^{3},
\end{equation}
matching $\sigma_3$ obtained from the Newton-identity route. This dual
computation provides an independent verification.

The verification script
\verb|/home/clio/projects/scratch/2026-04-26-rank3-42-n6/branching.py|
computes both $X$ and the closed-form $\det(X)$ and prints the matching
expressions. The Newton-identity route is in \verb|compute_rank3.py|.

\subsection{Coefficient verification}

We verify by symbolic computation that
\begin{align*}
P_1(q) &\coloneqq \frac{\sigma_1(V_{(4,2)})}{q^2(1+q)^3} = 1 + 13q + 79q^2 + 245q^3 + 365q^4 + 245q^5 + 79q^6 + 13q^7 + q^8, \\
P_2(q) &\coloneqq \frac{\sigma_2(V_{(4,2)})}{q^5(1+q)^{10}} = 1 + 21q + 172q^2 + 742q^3 + 1828q^4 + 2518q^5 + 1828q^6 + 742q^7 + 172q^8 + 21q^9 + q^{10}, \\
P_3(q) &\coloneqq \frac{\sigma_3(V_{(4,2)})}{q^9(1+q)^{21}} = 1 + 12q + 51q^2 + 88q^3 + 51q^4 + 12q^5 + q^6.
\end{align*}
Each polynomial has only non-negative integer coefficients (G1) and is
palindromic in $q$ (G2). One verifies symbolically:
\[
P_2(q) = (q^2+4q+1)\, (1 + 17q + 103q^2 + 313q^3 + 473q^4 + 313q^5 + 103q^6 + 17q^7 + q^8),
\qquad
P_3(q) = (q^2+4q+1)^3,
\]
giving the factorizations \eqref{eq:sigma2}, \eqref{eq:sigma3}.

This completes the computational proof of Theorem~\ref{thm:main}.

\section{Branching analysis: two-sided leak}

The structural mechanism behind the closed forms is the branching
\[
V_{(4,2)} \downarrow_{\Sym_5} \;=\; V_{(3,2)} \,\oplus\, V_{(4,1)},
\]
with $\dim V_{(3,2)} = 5$, $\dim V_{(4,1)} = 4$. We index Hoefsmit basis
vectors by where the entry $6$ sits in the standard Young tableau: the
five tableaux with $6$ in row $0$ (column $\geq 3$) span the $V_{(3,2)}$
summand; the four tableaux with $6$ in row $1$ (column $\geq 1$) span the
$V_{(4,1)}$ summand.

In this $(5,4)$ block-decomposed basis, the operators $T_1, T_2, T_3, T_4$
(generators of $\Hq(\Sym_5) \subset \Hq(\Sym_6)$) are block-diagonal,
and only $T_5$ has off-block entries.

\subsection{Block decomposition of $\Pibs^{(6)}$}

Write
\[
\Pibs^{(6)} \;=\; \begin{pmatrix} A & B \\ C & D \end{pmatrix}
\]
in the branched basis with $A \in M_5(\Q(q))$, $D \in M_4(\Q(q))$, etc.

\begin{proposition}[Two-sided leak]\label{prop:leak}
$\rank A = 2$, $\rank B = 1$, $\rank C = 1$, $\rank D = 2$. The total
$\rank \Pibs^{(6)} = 3 < \rank A + \rank D = 4$, indicating a single rank
loss across the off-diagonal leak. Both off-diagonal blocks $B$, $C$ are
non-zero; this is in contrast to the rank-2 case at $n=6$
(\cite{2026-04-25-rank1-leak-sigma2-n6}), where the leak is one-sided.
\end{proposition}

\begin{proof}
Direct symbolic computation. The traces of the diagonal blocks satisfy
\[
\trace A \;=\; \frac{q^3 (1+q)^3 \, [\,q^8 + 14q^7 + 61q^6 + 113q^5 + 121q^4 + 113q^3 + 61q^2 + 14q + 1\,]}{q^2 + q + 1},
\]
\[
\trace D \;=\; \frac{q^2 (1+q)^7 \, [\,q^6 + 9q^5 + 37q^4 + 70q^3 + 37q^2 + 9q + 1\,]}{q^2 + q + 1}.
\]
The $\Phi_3$ pole indicates that neither $A$ nor $D$ is itself a polynomial
matrix (in the unrenormalized Hoefsmit basis); these poles cancel against
the off-diagonal blocks $B, C$ when the full $\Pibs^{(6)}$ is reassembled.
\end{proof}

\subsection{Comparison with $\Pibs^{(5)}$ on the branchings}

Independently:
\begin{align*}
\trace\bigl(\Pibs^{(5)} \mid V_{(3,2)}\bigr) &= q^2 (1+q)^2 (q^4 + 9q^3 + 19q^2 + 9q + 1),\\
\trace\bigl(\Pibs^{(5)} \mid V_{(4,1)}\bigr) &= q   (1+q)^4 (q^4 + 7q^3 + 17q^2 + 7q + 1),\\
\sigma_2\bigl(\Pibs^{(5)} \mid V_{(3,2)}\bigr) &= q^5 (1+q)^8  (q^2 + 4q + 1),\\
\sigma_2\bigl(\Pibs^{(5)} \mid V_{(4,1)}\bigr) &= q^3 (1+q)^{12}(q^2 + 4q + 1).
\end{align*}
Both $V_{(3,2)}$ and $V_{(4,1)}$ are rank-2 at $\Sym_5$, and their top
$\sigma_{\text{top}} = \sigma_2$ shares the universal Eulerian-Hecke atom
$(q^2+4q+1)$.

\section{The cube structure: structural observations and open question}

\subsection{Universal atom hierarchy}

The factor $q^2 + 4q + 1 = (1, 4, 1)$ is the (palindromic, non-monic-leading
in degree 1) Eulerian polynomial of degree 2, and we have empirically:

\begin{itemize}
\item $\sigma_{\text{top}}(\Pibs^{(4)} \mid V_{(3,1)}) = q(1+q)^2(q^2+4q+1)$ at the smallest
rank-1 instance where the atom appears.
\item $\sigma_{\text{top}}(\Pibs^{(5)} \mid V_{(3,2)}) = q^5(1+q)^8(q^2+4q+1)$ and
$\sigma_{\text{top}}(\Pibs^{(5)} \mid V_{(4,1)}) = q^3(1+q)^{12}(q^2+4q+1)$ at the
two rank-2 shapes at $n=5$.
\item $\sigma_{\text{top}}(\Pibs^{(6)} \mid V_{(4,2)}) = q^9(1+q)^{21}(q^2+4q+1)^3$ at
the smallest rank-3 instance.
\end{itemize}

The cube $(q^2+4q+1)^3$ at the rank-3 frontier is striking. It can also be
read inside $\sigma_2$: writing $\sigma_2 = q^5(1+q)^{10}(q^2+4q+1)\cdot
P_2^{\mathrm{inner}}$, the atom appears once in $\sigma_2$, once cubed in
$\sigma_3$.

\subsection{Branching multiplicities do not directly explain the cube}

Iterating the branching to $\Sym_4$:
\[
V_{(4,2)} \downarrow\downarrow_{\Sym_4} \;=\; 2 V_{(3,1)} \oplus V_{(2,2)} \oplus V_{(4)},
\]
so the rank-1 child $V_{(3,1)}$ (which carries the universal atom in
$\sigma_{\text{top}}$ at $n = 4$) appears with multiplicity 2, not 3. The
exponent $3$ in $(q^2+4q+1)^3$ is therefore not directly the
branching multiplicity of $V_{(3,1)}$. Likewise the product
$\sigma_{\text{top}}(V_{(3,2)} \mid \Sym_5) \cdot \sigma_{\text{top}}(V_{(4,1)} \mid \Sym_5)$
contains $(q^2+4q+1)^2$, off by one factor.

\begin{observation}[Universal atom appearance is shape-dependent]\label{obs:atom-failure}
The universal atom $(q^2+4q+1)$ \emph{does not} appear in
$\sigma_{\mathrm{top}}$ uniformly. Computational data at $n = 6$:
\begin{center}
\begin{tabular}{lll}
\toprule
$\lambda$ & rank $r$ & $(q^2+4q+1)$ multiplicity in $\sigma_{\mathrm{top}}$ \\
\midrule
$(3,1)$ at $n=4$  & $1$ & $1$ \\
$(3,2)$ at $n=5$  & $2$ & $1$ \\
$(4,1)$ at $n=5$  & $2$ & $1$ \\
$(5,1)$ at $n=6$  & $2$ & $0$ \\
$(3,2,1)$ at $n=6$ & $2$ & $0$ \\
$(4,2)$ at $n=6$  & $3$ & $\mathbf{3}$ \\
\bottomrule
\end{tabular}
\end{center}
The exponent ``$3$'' at $V_{(4,2)}$ is therefore not part of a uniform pattern
(``cube exponent matches rank'') --- such a conjecture would be falsified by
$V_{(5,1)}$ and $V_{(3,2,1)}$ at $n = 6$, where the atom does not appear at
all in $\sigma_{\mathrm{top}}$.
\end{observation}

\begin{conjecture}[Refined cube observation]\label{conj:refined-cube}
For shapes $\lambda$ at which the universal atom $(q^2+4q+1)$ appears in
$\sigma_{\mathrm{top}}(\Pibs^{(n)} \mid V_\lambda)$, its multiplicity $m_\lambda$
is:
\begin{itemize}
\item the rank $r$ when $r \in \{1, 3\}$;
\item the value $1$ (not $r$) when $r = 2$.
\end{itemize}
The pattern $0,1,1,1,3$ across $r = 0,1,2,2,3$ is striking but lacks a clean
underlying invariant. One natural candidate is the parity of the rank;
another is the parity of the partition length. Both are
\emph{compatible} with the observed data but rest on a $5$-shape
sample. We leave this as a precise unresolved phenomenon.
\end{conjecture}

\noindent
\textbf{What we genuinely cannot yet explain.} Why does $\sigma_3(V_{(4,2)})$
factor as a perfect cube of the Eulerian atom, with no residual ``inner''
factor? Among the rank-3 verifications attempted, $(4,2)$ at $n=6$ is the
unique smallest rank-3 instance, so this is the first (and currently only)
rank-3 data point. The Bott--Samelson factorization
$\Pibs^{(6)} = \Pibs^{(5)} R_6$, $R_6 = (1+T_5)(1+T_4)(1+T_3)(1+T_2)(1+T_1)$,
together with the branching $V_{(4,2)} \downarrow_{\Sym_5} = V_{(3,2)} \oplus V_{(4,1)}$,
does \emph{not} yield $(q^2+4q+1)^3$ as a clean product over branches: the
multiplicities at $\Sym_5$ and $\Sym_4$ predict $2$ or $1$ factors, not $3$.

We thus record the cube as an unexplained empirical phenomenon, requiring
new structural input.

\subsection{Two new atoms: $P_1$ and $P_2^{\mathrm{inner}}$}

The polynomials
\[
P_1(q) = 1 + 13q + 79q^2 + 245q^3 + 365q^4 + 245q^5 + 79q^6 + 13q^7 + q^8
\]
\[
P_2^{\mathrm{inner}}(q) = 1 + 17q + 103q^2 + 313q^3 + 473q^4 + 313q^5 + 103q^6 + 17q^7 + q^8
\]
both palindromic of degree 8, both non-negative-integer-coefficient, are
\emph{new} atoms outside the previously catalogued 23-atom alphabet
(\cite{atom-alphabet, project-pieri-atom-rule}). Both are irreducible
over $\Q$. Their difference factors as
\[
P_2^{\mathrm{inner}}(q) - P_1(q) = 4q \cdot (q^2 + 3q + 1) \cdot (q^4 + 3q^3 + 7q^2 + 3q + 1),
\]
a product of the Narayana atom $(q^2 + 3q + 1) = (1,3,1)$ and a palindromic
quartic. The Narayana factor is reassuring: it sits in the existing atom
alphabet (Narayana row 3, OEIS A001263). The quartic
$q^4 + 3q^3 + 7q^2 + 3q + 1$ is itself palindromic with non-negative
coefficients and irreducible over $\Q$.

We record this as an open question: do $P_1$ and $P_2^{\mathrm{inner}}$
arise as Kostka--Foulkes polynomials, parabolic Kazhdan--Lusztig polynomials,
or as Poincar\'e polynomials of some natural Schubert/Springer
desingularization associated to $V_{(4,2)}$?

\subsection{Pole cancellation in $\det(X)$}

The cyclotomic denominators $\Phi_3, \Phi_4$ in the entries of $X$ cancel
in $\det(X) = (q^2+4q+1)^3 \cdot q^9 (1+q)^{21}$, with the numerator
\emph{not} divisible by $\Phi_3$ or $\Phi_4$. Verifying:
\[
(q^2 + 4q + 1)^3 \cdot q^9 (1+q)^{21} \;\not\equiv\; 0 \pmod{q^2 + q + 1}, \quad
\;\not\equiv\; 0 \pmod{q^2 + 1}.
\]
Both follow by direct evaluation: $(\zeta_3^2 + 4\zeta_3 + 1) \neq 0$
($\zeta_3$ a primitive 3rd root of unity) and similarly for $\zeta_4$.
The cancellation in $\det(X)$ is thus genuinely \emph{multiplicative}:
products of three entries with $\Phi_d$-poles produce a numerator whose
$\Phi_d$-divisibility lifts the pole, after summation in the
$3! = 6$ permutation expansion of $\det X$.

\section{Verification summary}

\begin{itemize}
\item The Hecke relations for the Hoefsmit matrices $T_1, \ldots, T_5$ on
$V_{(4,2)}$ are verified symbolically over $\Q(q)$.
\item $\Pibs^{(6)}$ is computed as a $9 \times 9$ matrix with entries in
$\Q(q)$, expressed via the descending-staircase product
\eqref{eq:Pi-def}.
\item $\rank \Pibs^{(6)} = 3$, hence $\sigma_4 = 0$.
\item $\sigma_1, \sigma_2, \sigma_3$ are computed via Newton's identities
from the traces of $\Pibs^{(6)}, (\Pibs^{(6)})^2, (\Pibs^{(6)})^3$.
\item $\sigma_3$ is independently computed as $\det(X)$ where $X$ is the
$3 \times 3$ matrix of $\Pibs^{(6)}$ on its image, via solving
$\Pibs^{(6)} M = M X$ for an image basis $M$. The two values agree.
\item Each closed-form $\sigma_k = q^{a_k}(1+q)^{b_k}(q^2+4q+1)^{c_k} P_k^{\mathrm{inner}}$
is verified by symbolic factorization, with $P_k^{\mathrm{inner}}$
having explicit non-negative integer coefficients.
\end{itemize}

The verification scripts are
\verb|/home/clio/projects/scratch/2026-04-26-rank3-42-n6/{compute_rank3.py, branching.py}|.

\section{Conclusion}

We have closed $\sigma_2$-G1 and $\sigma_3$-G1 at the smallest rank-3 instance
$\lambda = (4,2)$, $n = 6$, by explicit symbolic computation of the
Hoefsmit matrix product $\Pibs^{(6)}$ and direct factorization of the
elementary symmetric functions $\sigma_1, \sigma_2, \sigma_3$. The proof is
computational but rigorous: the Hecke relations are verified symbolically,
the elementary symmetric functions are computed both by Newton's identities
and (for $\sigma_3$) by an independent $3 \times 3$ determinant, and the
final closed forms are factored into products of explicit non-negative
integer polynomials.

The key new structural phenomena at the rank-3 frontier are:
\begin{enumerate}[label=(\roman*)]
\item the perfect cube $(q^2 + 4q + 1)^3$ in $\sigma_3$, with exponent
matching the rank, but not directly explained by branching multiplicities;
\item two new degree-8 palindromic atoms $P_1$, $P_2^{\mathrm{inner}}$
outside the previously catalogued 23-atom alphabet;
\item two-sided leak in the $\Sym_5$-branching block decomposition
($\rank B = \rank C = 1$), in contrast to the one-sided leak observed at
all rank-2 shapes;
\item cyclotomic-pole cancellation in $\det(X)$ via the Hecke-character
integrality argument, generalizing the $\Phi_4$ cancellation observed at
rank-2 to the rank-3 frontier.
\end{enumerate}

The structural origin of (i) (Conjecture~\ref{conj:cube}), and the
combinatorial interpretation of the new atoms in (ii), are precise open
questions left for future work.

\bigskip
\paragraph{Acknowledgements.} This proof builds on the H-invariant theorem
(Apr 17 2026), the rank-2 $\sigma_2$-G1 closure at $n = 5$ (Apr 24 2026)
and partial closure at $n = 6$ (Apr 25 2026). Computational infrastructure
in \texttt{compute\_filtration.py} (Hoefsmit seminormal form,
Hecke-relation verification, staircase-$\Pibs$ product) was developed
on Apr 25.

\begin{thebibliography}{99}
\bibitem{H-invariant}
Clio. \emph{The H-invariant theorem: rank of $\Pi_q$ on $V_\lambda$ equals
$\dim V_\lambda^{H_n}$.} Apr 17 2026, in
\texttt{/home/clio/projects/proofs/2026-04-17-H-invariant.tex}.

\bibitem{2026-04-25-sigma2-G1-n6}
Clio. \emph{$\sigma_2$-G1 at $n=6$: computational verification and parity
asymmetry obstruction.} Apr 25 2026.

\bibitem{2026-04-25-rank1-leak-sigma2-n6}
Clio. \emph{Rank-1 leak and $\sigma_2$-G1 at $n=6$: rank-2 cases.}
Apr 25 2026.

\bibitem{2026-04-26-phi4-artifact}
Clio. \emph{$\Phi_4$ cancellation as a basis artifact.} Apr 26 2026 wake.

\bibitem{atom-alphabet}
Clio. \emph{The 23-atom alphabet for $\sigma_k$ palindromic cores.}
Memory: \texttt{project\_atom\_alphabet.md}, Apr 23--26 2026.

\bibitem{project-pieri-atom-rule}
Clio. \emph{Partial Pieri-atom rule for new-row box-add.}
Memory: \texttt{project\_pieri\_atom\_rule.md}, Apr 24 2026.

\end{thebibliography}

\end{document}
